\chapter{La tension et le voltmètre}\label{ch:g8:voltage-voltmeter}

Depuis que la première \omterm{def:g3:first-electric-circuit:battery}{pile} a allumé la première lampe, un mot fait
dans ce livre un service lourd et vague~: la «~poussée~». Les \omterm{def:g3:first-electric-circuit:battery}{piles}
poussent faiblement, le secteur pousse mortellement, deux \omterm{def:g3:first-electric-circuit:battery}{piles} bout à
bout poussent ensemble. La poussée devient aujourd'hui une grandeur
mesurée, munie d'une \omterm{def:g6:measurement-in-science:measuring}{unité} célèbre --- et son instrument, contrairement
à celui du chapitre précédent, refuse de se tenir sur la route.

\section{La tension}

\begin{definition}[Tension]\label{def:g8:voltage-voltmeter:voltage}
La \emph{tension}\index{tension} $U$ entre deux points d'un circuit
mesure la poussée électrique disponible entre eux --- avec quelle
\omterm{def:g2:pushes-and-pulls:force}{force} le courant est chassé d'un point vers l'autre, comme la hauteur
d'une chute d'eau mesure avec quelle \omterm{def:g2:pushes-and-pulls:force}{force} elle entraîne la roue du
moulin. Une tension est toujours un \emph{entre}~: entre les deux
\omterm{def:g3:first-electric-circuit:battery}{bornes} d'une \omterm{def:g3:first-electric-circuit:battery}{pile}, entre les deux côtés d'une lampe --- jamais en un
seul point. Son \omterm{def:g6:measurement-in-science:measuring}{unité} est le \emph{volt}\index{volt} (\unit{V}), en
l'honneur de l'inventeur de la première \omterm{def:g3:first-electric-circuit:battery}{pile}.
\end{definition}

\begin{example}[L'album de famille des volts]\label{ex:g8:voltage-voltmeter:values}
Une \omterm{def:g3:first-electric-circuit:battery}{pile} ronde~: \qty{1.5}{V}. La \omterm{def:g3:first-electric-circuit:battery}{pile} plate~: \qty{4.5}{V} --- trois
éléments bout à bout sous un même emballage. La \omterm{def:g3:first-electric-circuit:battery}{pile} rectangulaire de
neuf \omterm{def:g8:voltage-voltmeter:voltage}{volts}, ses deux \omterm{def:g3:first-electric-circuit:battery}{bornes} sur le dessus~: \qty{9}{V}, six éléments
cachés. Une batterie de voiture~: \qty{12}{V}. La prise murale~:
\qty{230}{V} --- plus de cent cinquante petits éléments de poussée, ce
qui est tout le contenu chiffré du «~la prise, jamais~». Un éclair
entre le nuage et le sol~: des centaines de millions de \omterm{def:g8:voltage-voltmeter:voltage}{volts}, un
bref instant.
\end{example}

\begin{proposition}[Une tension sans courant]\label{prop:g8:voltage-voltmeter:nocurrent}
La \omterm{def:g8:voltage-voltmeter:voltage}{tension} et le courant sont deux grandeurs différentes, et la
\omterm{def:g8:voltage-voltmeter:voltage}{tension} vient la première. Une \omterm{def:g3:first-electric-circuit:battery}{pile} seule dans un tiroir, reliée à
rien, ne fait circuler aucun courant --- et pourtant, entre ses
\omterm{def:g3:first-electric-circuit:battery}{bornes}, les \qty{1.5}{V} au complet se tiennent prêts, mesurables,
comme la hauteur d'une chute d'eau un jour où le moulin est fermé. La
poussée peut exister sans marche~; la marche n'existe jamais sans
poussée.
\end{proposition}

\begin{definition}[Voltmètre]\label{def:g8:voltage-voltmeter:voltmeter}
Un \emph{voltmètre}\index{voltmètre} mesure la \omterm{def:g8:voltage-voltmeter:voltage}{tension} entre deux
points --- il se \omterm{def:g5:series-parallel-circuits:parallel}{branche} donc \emph{à leurs \omterm{def:g3:first-electric-circuit:battery}{bornes}}~: \emph{\omterm{def:g5:series-parallel-circuits:parallel}{en
dérivation}} avec le composant étudié, un fil de chaque côté, comme un
géomètre tendant son niveau entre deux hauteurs. Construit à l'opposé
de l'\omterm{def:g8:intensity-ammeter:ammeter}{ampèremètre} en tout point, il s'oppose au courant presque
\emph{complètement}~: branché aux \omterm{def:g3:first-electric-circuit:battery}{bornes} de n'importe quoi, il n'en
prélève quasiment rien --- si bien que le poser ne dérange rien, et
qu'aucune boucle n'a jamais besoin d'être ouverte.
\end{definition}

\begin{omfigure}
\begin{tikzpicture}[scale=0.9]
  \draw (0,0) to[battery1, l=pile] (0,2.8)
    to[short] (2.0,2.8)
    to[lamp, l=lampe] (4.6,2.8)
    to[short] (6.4,2.8)
    to[short] (6.4,0)
    to[short] (0,0);
  \draw (2.0,2.8) to[short] (2.0,4.6)
    to[voltmeter, l=$U$] (4.6,4.6)
    to[short] (4.6,2.8);
\end{tikzpicture}

{\small Le \omterm{def:g8:voltage-voltmeter:voltmeter}{voltmètre} prend la posture du géomètre~: aux \omterm{def:g3:first-electric-circuit:battery}{bornes} de la
lampe, un fil de chaque côté, mesurant la poussée entre les deux ---
et jamais sur la route.}
\end{omfigure}

\begin{method}[Mesurer une tension]\label{met:g8:voltage-voltmeter:measure}
\begin{enumerate}
  \item laisser le circuit intact --- on n'ouvre rien pour un
        \omterm{def:g8:voltage-voltmeter:voltmeter}{voltmètre}~;
  \item poser un fil de chaque côté du composant~: le fil \unit{V}
        là où le courant conventionnel entre, le fil commun là où
        il sort~;
  \item plus grand calibre d'abord, puis descendre jusqu'à la
        lecture confortable~;
  \item lire, et écrire l'\omterm{def:g6:measurement-in-science:measuring}{unité}.
\end{enumerate}
Les manières des deux appareils, côte à côte~: l'\omterm{def:g8:intensity-ammeter:ammeter}{ampèremètre} ---
ouvrir la route, s'y tenir, ne s'opposer à rien~; le \omterm{def:g8:voltage-voltmeter:voltmeter}{voltmètre} ---
n'ouvrir rien, se tenir à côté, s'opposer à tout. Confondre les deux
manières est l'habitude la plus coûteuse de l'atelier.
\end{method}

\section{Les tensions nominales}

\begin{definition}[Tension nominale]\label{def:g8:voltage-voltmeter:nominal}
Toute lampe, tout \omterm{ex:g6:circuits-and-safety:motor}{moteur}, tout appareil porte une \emph{tension
nominale}\index{tension nominale} --- la poussée pour laquelle il est
construit, gravée sur son culot ou sa plaque~: une ampoule de lampe
de poche «~3,5 V~», un phare de voiture «~12 V~», une bouilloire
«~230 V~». Alimenté à sa \omterm{def:g8:voltage-voltmeter:voltage}{tension} nominale, l'appareil se comporte
comme prévu~; alimenté à bien moins, il boude --- lampe terne, \omterm{ex:g6:circuits-and-safety:motor}{moteur}
paresseux~; alimenté à bien plus, il grille --- brillamment, une
fois.
\end{definition}

\begin{example}[Accorder la lampe et la pile]\label{ex:g8:voltage-voltmeter:matching}
Une ampoule «~3,5 V~» sur la \omterm{def:g3:first-electric-circuit:battery}{pile} plate de \qty{4.5}{V}~: un peu
surmenée --- brillante, et de courte vie. Sur un seul élément de
\qty{1.5}{V}~: une lueur orange endormie. Sur la \omterm{def:g3:first-electric-circuit:battery}{pile} de
\qty{9}{V}~: un instant éblouissant, puis le fil ténu se rompt ---
la nécrologie du filament. L'étiquette est un contrat~: honore-le, et
la lampe vivra la vie qu'elle annonce.
\end{example}

\begin{remark}[Pourquoi 230 V tue]\label{rem:g8:voltage-voltmeter:mains}
La vieille règle absolue reçoit enfin son nombre. La route
conductrice humide du corps, offerte à \qty{4.5}{V}, laisse passer un
courant trop faible pour être senti. La même route sous \qty{230}{V}
--- cinquante fois la poussée --- laisse passer cinquante fois le
courant, bien au-delà des quelques centièmes d'\omterm{def:g8:intensity-ammeter:intensity}{ampère} qui arrêtent un
cœur. La prise n'est pas «~un peu~» plus forte que ton établi à
\omterm{def:g3:first-electric-circuit:battery}{piles}~: elle est deux ordres de grandeur au-dessus, et la sévérité du
règlement n'est que de l'arithmétique.
\end{remark}

\section{Les premières lectures de tension}

\begin{example}[Le tour d'une boucle simple]\label{ex:g8:voltage-voltmeter:readings}
Une \omterm{def:g3:first-electric-circuit:battery}{pile}, une lampe, les appareils en main. Aux \omterm{def:g3:first-electric-circuit:battery}{bornes} de la \omterm{def:g3:first-electric-circuit:battery}{pile}~:
\qty{4.5}{V}. Aux \omterm{def:g3:first-electric-circuit:battery}{bornes} de la lampe~: \qty{4.5}{V} --- la lampe
reçoit ce que la \omterm{def:g3:first-electric-circuit:battery}{pile} offre. Aux \omterm{def:g3:first-electric-circuit:battery}{bornes} d'un tronçon de \emph{fil} de
liaison~: \qty{0}{V}, ou aussi près que l'appareil sache le dire ---
un bon fil est une route plate, qui ne dépense rien de la poussée. La
poussée vit aux \omterm{def:g3:first-electric-circuit:battery}{bornes} des \emph{travailleurs}, non le long des
routes~: une lecture à retenir pour le moment où les lois du chapitre
suivant la rendront officielle.
\end{example}

\begin{example}[L'aveu de l'interrupteur ouvert]\label{ex:g8:voltage-voltmeter:openswitch}
Lampe éteinte, \omterm{ex:g3:first-electric-circuit:switch}{interrupteur} ouvert --- où est passée la poussée~?
Mesurons~: aux \omterm{def:g3:first-electric-circuit:battery}{bornes} de la lampe éteinte, \qty{0}{V}~; aux \omterm{def:g3:first-electric-circuit:battery}{bornes} de
l'\emph{\omterm{ex:g3:first-electric-circuit:switch}{interrupteur} ouvert}, les \qty{4.5}{V} au complet. La poussée
de la \omterm{def:g3:first-electric-circuit:battery}{pile} se tient tout entière en travers de la coupure, à
attendre --- le seul endroit de la boucle où deux points soient
vraiment «~éloignés~» électriquement. Ferme l'\omterm{ex:g3:first-electric-circuit:switch}{interrupteur}, et les
deux lectures échangent leurs places. Un \omterm{ex:g3:first-electric-circuit:switch}{interrupteur} ouvert est
l'endroit où la \omterm{def:g8:voltage-voltmeter:voltage}{tension} attend.
\end{example}

\section{Exercices}

\begin{exercise}[$\star$]\label{exo:g8:voltage-voltmeter:1}
Que mesure une \omterm{def:g8:voltage-voltmeter:voltage}{tension}, et pourquoi est-elle toujours «~entre~»~?
Donne son \omterm{def:g6:measurement-in-science:measuring}{unité} et l'\omterm{prop:g5:mirrors-reflection:image}{image} de la chute d'eau.
\end{exercise}

\begin{exercise}[$\star$]\label{exo:g8:voltage-voltmeter:2}
De mémoire~: les \omterm{def:g8:voltage-voltmeter:voltage}{tensions} d'une \omterm{def:g3:first-electric-circuit:battery}{pile} ronde, de la \omterm{def:g3:first-electric-circuit:battery}{pile} plate, de la
\omterm{def:g3:first-electric-circuit:battery}{pile} rectangulaire, d'une batterie de voiture, du secteur.
\end{exercise}

\begin{exercise}[$\star$]\label{exo:g8:voltage-voltmeter:3}
Une \omterm{def:g3:first-electric-circuit:battery}{pile} est posée là, reliée à rien. Quel courant circule~? Que lit
un \omterm{def:g8:voltage-voltmeter:voltmeter}{voltmètre} à ses \omterm{def:g3:first-electric-circuit:battery}{bornes}~? Quelle proposition prend la parole~?
\end{exercise}

\begin{exercise}[$\star$]\label{exo:g8:voltage-voltmeter:4}
Oppose les manières des deux appareils~: branchement, ouverture de la
route, opposition au courant. Quelle erreur fabrique un
\omterm{def:g7:short-circuits-safety:device}{court-circuit} muni d'un cadran~?
\end{exercise}

\begin{exercise}[$\star$]\label{exo:g8:voltage-voltmeter:5}
Qu'est-ce qu'une \omterm{def:g8:voltage-voltmeter:nominal}{tension nominale}~? Prédis le comportement d'une
lampe «~6 V~» sous \qty{1.5}{V}~; sous \qty{6}{V}~; sous
\qty{12}{V}.
\end{exercise}

\begin{exercise}[$\star$]\label{exo:g8:voltage-voltmeter:6}
Dans la boucle simple, que lit le \omterm{def:g8:voltage-voltmeter:voltmeter}{voltmètre} aux \omterm{def:g3:first-electric-circuit:battery}{bornes} de la \omterm{def:g3:first-electric-circuit:battery}{pile}, de
la lampe, et d'un tronçon de fil~? Où vit la poussée~?
\end{exercise}

\begin{exercise}[$\star$]\label{exo:g8:voltage-voltmeter:7}
\omterm{ex:g3:first-electric-circuit:switch}{Interrupteur} ouvert~: que trouve-t-on aux \omterm{def:g3:first-electric-circuit:battery}{bornes} de la lampe éteinte,
et aux \omterm{def:g3:first-electric-circuit:battery}{bornes} de l'\omterm{ex:g3:first-electric-circuit:switch}{interrupteur} ouvert~? Qu'arrive-t-il aux deux
lectures au moment du déclic~?
\end{exercise}

\begin{exercise}[$\star$]\label{exo:g8:voltage-voltmeter:8}
Pourquoi poser un \omterm{def:g8:voltage-voltmeter:voltmeter}{voltmètre} aux \omterm{def:g3:first-electric-circuit:battery}{bornes} d'une lampe allumée ne
dérange-t-il presque pas le circuit~? Que refuse de faire cet
appareil~?
\end{exercise}

\begin{exercise}[$\star\star$]\label{exo:g8:voltage-voltmeter:9}
La \omterm{def:g3:first-electric-circuit:battery}{pile} plate de \qty{4.5}{V} cache trois éléments. Dessine (ou
décris) leur disposition, et calcule la \omterm{def:g8:voltage-voltmeter:voltage}{tension} si l'un des éléments
était retourné par accident --- le chapitre de la lampe de poche,
il y a des années, maintenant chiffré. (Des poussées opposées se
retranchent.)
\end{exercise}

\begin{exercise}[$\star\star$]\label{exo:g8:voltage-voltmeter:10}
La lampe d'un campeur porte la mention \qty{12}{V}~; la batterie du
fourgon indique \qty{12.6}{V} au départ et \qty{11.2}{V} après une
longue soirée. Décris la soirée de la lampe, et pourquoi rien n'a
grillé.
\end{exercise}

\begin{exercise}[$\star\star$]\label{exo:g8:voltage-voltmeter:11}
Un \omterm{def:g8:voltage-voltmeter:voltmeter}{voltmètre} branché par erreur \emph{\omterm{def:g4:series-circuits:series}{en série}} avec la lampe (boucle
ouverte, appareil dans la coupure)~: prédis le comportement de la
lampe et la lecture de l'appareil --- souviens-toi de ce à quoi
s'oppose un \omterm{def:g8:voltage-voltmeter:voltmeter}{voltmètre}. (Le silence de la lampe et le nombre de
l'appareil sont l'un et l'autre instructifs.)
\end{exercise}

\begin{exercise}[$\star\star\star$]\label{exo:g8:voltage-voltmeter:12}
Deux lampes identiques \omterm{def:g4:series-circuits:series}{en série} sous \qty{4.5}{V}~: un \omterm{def:g8:voltage-voltmeter:voltmeter}{voltmètre}
trouve environ \qty{2.25}{V} aux \omterm{def:g3:first-electric-circuit:battery}{bornes} de chacune. Les mêmes lampes
\omterm{def:g5:series-parallel-circuits:parallel}{en dérivation} sur la même \omterm{def:g3:first-electric-circuit:battery}{pile}~: prédis la \omterm{def:g8:voltage-voltmeter:voltage}{tension} de chaque lampe,
et concilie les deux constats avec les règles d'éclat que tu connais
depuis des années (une paire \omterm{def:g4:series-circuits:series}{en série} s'assombrit, une paire \omterm{def:g5:series-parallel-circuits:parallel}{en
dérivation} brille à plein). Quel chapitre à venir fera de tes
prédictions une loi~?
\end{exercise}

\section{Problème~: l'audit du tiroir à piles}

\begin{problem}[{Problème du week-end --- le grand audit du tiroir à
piles~; un voltmètre, une boîte à chaussures pleine d'éléments et
l'étagère à gadgets de la famille}]\label{pb:g8:voltage-voltmeter:1}
Toute maison a ce tiroir~: des éléments en vrac, des \omterm{def:g3:first-electric-circuit:battery}{piles}
mystérieuses, des gadgets qui «~ont besoin de neuves~». Jour d'audit
--- \omterm{def:g8:voltage-voltmeter:voltmeter}{voltmètre} en main.

\medskip\noindent\textbf{Partie I --- trier les éléments.}
\begin{enumerate}
  \item Décris le branchement correct du \omterm{def:g8:voltage-voltmeter:voltmeter}{voltmètre} pour contrôler un
        élément de \qty{1.5}{V} en vrac --- et pourquoi il n'y a
        aucun circuit à monter d'abord.
  \item Les lectures de l'audit, sorties du tiroir~: \qty{1.58}{V},
        \qty{1.32}{V}, \qty{0.9}{V}, \qty{0.1}{V}. Classe-les~:
        neuve, utilisable, fatiguée, morte.
  \item Une \omterm{def:g3:first-electric-circuit:battery}{pile} de neuf \omterm{def:g8:voltage-voltmeter:voltage}{volts} indique \qty{8.9}{V}. Que fournit à
        peu près chacun de ses six éléments intérieurs~?
  \item Pourquoi un élément presque mort montre-t-il encore
        \emph{une} \omterm{def:g8:voltage-voltmeter:voltage}{tension} alors qu'il ne peut plus rien allumer~?
        (Sa poussée survit faiblement~; qu'a-t-il perdu la \omterm{def:g2:pushes-and-pulls:force}{force}
        d'entraîner~?)
\end{enumerate}

\medskip\noindent\textbf{Partie II --- l'étagère à gadgets.}
\begin{enumerate}[resume]
  \item L'appareil photo veut deux éléments bout à bout. Quelle
        \omterm{def:g8:voltage-voltmeter:voltage}{tension} attend-il, et que délivrera la paire de rescapés à
        \qty{1.32}{V} de la question 2~?
  \item Le vieux poste de radio prend quatre éléments bout à bout et
        boude en dessous de \qty{5}{V}. Quels éléments du tiroir
        peuvent, s'il y en a, former son équipage~?
  \item Le petit frère tasse trois éléments dans la lampe de poche
        qui en prévoit deux, «~pour plus de \omterm{def:g1:light-and-shadows:source}{lumière}~», en coinçant
        le troisième contre le ressort. L'ampoule est donnée pour
        \qty{3}{V}. Prédis le spectacle, en langage de \omterm{def:g8:voltage-voltmeter:nominal}{tension
        nominale}.
  \item La plaque du transformateur de sonnette indique
        «~\qty{8}{V}~». Quelqu'un propose de brancher la sonnette
        directement sur une \omterm{def:g3:first-electric-circuit:battery}{pile} de neuf \omterm{def:g8:voltage-voltmeter:voltage}{volts}, «~puisque huit,
        c'est presque neuf~». Juge~: accord honnête ou meurtre
        lent~?
\end{enumerate}

\medskip\noindent\textbf{Partie III --- les lois de l'audit.}
\begin{enumerate}[resume]
  \item Dans la lampe de poche (ampoule et éléments \omterm{def:g4:series-circuits:series}{en série},
        \omterm{ex:g3:first-electric-circuit:switch}{interrupteur} ouvert), l'auditeur mesure~: aux \omterm{def:g3:first-electric-circuit:battery}{bornes} de
        l'\omterm{ex:g3:first-electric-circuit:switch}{interrupteur} ouvert, \qty{3.1}{V}~; aux \omterm{def:g3:first-electric-circuit:battery}{bornes} de
        l'ampoule, \qty{0}{V}. Explique les deux, et prédis les
        deux après le déclic.
  \item Aux \omterm{def:g3:first-electric-circuit:battery}{bornes} de chaque tronçon du câblage de la lampe de
        poche, l'appareil dit \qty{0}{V}, allumée comme éteinte. Que
        déclare cela sur un bon fil, en langage de poussée~?
  \item La prise du secteur n'est \emph{pas} auditée, par loi
        familiale permanente. Redis cette loi avec le nombre de ce
        chapitre, et la phrase arithmétique qui la fonde.
  \item Dépose le rapport d'audit~: trois phrases --- ce qu'un
        \omterm{def:g8:voltage-voltmeter:voltmeter}{voltmètre} est et n'est jamais~; ce qu'une \omterm{def:g8:voltage-voltmeter:nominal}{tension nominale}
        d'étiquette promet et menace~; et où, dans tout circuit au
        repos, la \omterm{def:g8:voltage-voltmeter:voltage}{tension} attend.
\end{enumerate}
\end{problem}
